\subsection{Checking the type prerequisites of C.20 and C.22}
The supplied \textsc{INDLOCAL} extract starts at C.20. Its references to
C.1--C.19 are therefore not proofs present in that extract. We supply
the precise prerequisites used below, and then check the written first
proof of C.20. This avoids treating an unavailable earlier argument as
an authority.

\begin{lemma}[The type-map contraction]\label{F:type-contraction}
A type is a finite word $\sigma$ in $+,-$, including the empty word.
Its interval $I_\sigma$ is the free category on the oriented line with
vertices $0,\ldots,|\sigma|$. Maps of types are endpoint-preserving,
nondecreasing functors of these intervals. For fixed $\sigma,\tau$, the
category $P(\sigma,\tau)$ of such functors and endpoint-fixed natural
transformations is either empty or categorically contractible: its
identity is connected to a constant functor by a finite zigzag of
natural transformations.
\end{lemma}
\begin{proof}
Write $n=|\sigma|$, $m=|\tau|$. A type map is equivalent to a weakly
increasing allocation
\[
 a_1\leq\cdots\leq a_m,\qquad
 a_j\in\{1,\ldots,n\},\quad \sigma_{a_j}=\tau_j.
\]
Target edge $j$ is traversed by source edge $a_j$. Conversely the image
of source vertex $i$ is the number of allocated target edges with
$a_j\leq i$. Each source edge then traverses a string of its own sign
or an identity, so this is exactly the required functor. This also
handles an empty target; a nonempty target cannot be allocated to an
empty source.

The interval category is a poset. A transformation from allocation
$a$ to allocation $b$ therefore exists uniquely precisely when
\[
 a_j\geq b_j\quad(\tau_j=+),\qquad
 a_j\leq b_j\quad(\tau_j=-).
\]
Indeed a component crosses edge $j$ in the positive vertex direction
exactly when its allocation has moved earlier, and can cross it in that
direction exactly when its sign is positive. The reverse statement
holds for negative edges. Pointwise existence implies naturality in a
poset, so these inequalities give both necessity and sufficiency.

Suppose allocations exist. Let $q_j$ be the minimum of their $j$th
coordinates. Taking these minima preserves monotonicity and the required
colors, so $q$ is an allocation.
Let $F_j(a)$ replace the first $j$ coordinates of $a$ by those of $q$.
The junction remains increasing, since $q_j\leq a_j\leq a_{j+1}$.
The displayed inequalities show that every $F_j$ is an order-preserving
endofunctor. For a positive $\tau_j$ there is
$F_{j-1}\Rightarrow F_j$; for a negative $\tau_j$ there is the reverse
transformation. Since $F_0=1$ and $F_m$ is constant at $q$, this is the
claimed contraction. No connectedness or coherence assertion is left
to a choice of individual paths.
\end{proof}

\begin{proposition}[The uniform fibred comparison used in C.28]
\label{F:C23-uniform}
For $u:J\to T$, the map of fibred Grothendieck constructions
\[
 \int_J^{\mathrm{fib}}F u^{\op}\longrightarrow
 \int_T^{\mathrm{fib}}F
\]
is a nerve weak equivalence for every strict $F:T^{\op}\to\Cat$
if and only if each comma $u/t$ is weakly contractible. Parallel
pullback homotopies and a cofinal reachability sequence alone do not
imply this condition. In particular they do not justify the universal
comparison proposed at VII ms.~99--110.
\end{proposition}
\begin{proof}
For $(t,z)$ in the target total category, forget fibre data in its
comma category. The resulting functor to $u/t$ has the section sending
$(j,\alpha:u(j)\to t)$ to the identity of $F(\alpha)z$.
The fibre data of any object are exactly an arrow
$x\to F(\alpha)z$. This arrow gives a natural transformation from
the identity to section after projection. A comma arrow $v$ obeys
$\alpha=\alpha'u(v)$, and contravariance gives the exact naturality
equation $F(\alpha)=F(u(v))F(\alpha')$. The comma spaces are therefore
equivalent to $B(u/t)$. Quillen's Theorem A proves sufficiency.
For necessity use the discrete presheaf $F(s)=\Hom_T(s,t)$: its
total category is $T/t$, with terminal object, and its pullback total
category is exactly $u/t$.

For the printed failure take $T=\mathbf BC_2$ and let
$u:\N^{\op}\to T$ send every arrow to the identity. Use the constant
terminal fibre. Parallel pullbacks agree, the reachability preorder
is a point, and a countable cofinal sequence exists. The source total
category is $\N^{\op}$ and is contractible, whereas the target is
$\mathbf BC_2$ with fundamental group $C_2$. Thus the specified map
is not a weak equivalence. For the separate smooth/proper suggestion,
the types $+-+$ and $+$ each have only identity endomorphisms, and
there are exactly two maps between them, choosing one of the two
positive edges. The corresponding incoming and outgoing lift
categories over their reachability arrow are discrete two-point
categories. They are not even connected. Lemma~\ref{F:type-contraction}
instead supplies transformations between type maps; it must not be
confused with either of these discrete lift categories.
\end{proof}

\begin{proposition}[C.22, the remaining type and colimit assertions]
\label{F:C22-types}
Concatenation of types and their interval functors is strictly
associative and unital. The empty type is terminal. The opposite type
category is neither filtered nor totally $W_\infty$-aspherical.
Nevertheless every actual projective sequence of types gives a
covariant path diagram indexed by the filtered category $\N$; its
ordinary categorical colimit has the homotopy type of its homotopy
colimit. Together with Proposition~\ref{F:C22}, these statements check
all assertions of \textup{\textsc{INDLOCAL}} C.22.
\end{proposition}
\begin{proof}
\textbf{Intervals and concatenation.} Gluing two free oriented lines at
their end and initial vertices gives their concatenated free line.
Functors out, and transformations between such functors, are exactly
pairs agreeing at the junction. Thus this is the asserted categorical
pushout and gives the corresponding pullback description of paths.
Blockwise concatenation of vertex maps is strictly associative, has the
empty block as strict unit, and commutes with composition. A type map
is uniquely a concatenation, over its source edges, of maps
$\varepsilon\to\varepsilon^l$, $l\geq0$. Under composition each such
string is refined by the lengths of the second map's strings.
Consequently the descriptions by functions $\{1,\ldots,l\}\to\N$
and their iterates in every nerve degree at VII ms.~50'--58 follow,
including zero lengths and identity insertions. Endpoint preservation
is exactly what preserves endpoint evaluations in an arrow
$(u,\eta):c\to d\circ u$ of the fibred path category.

\textbf{The restriction on interval maps.} The functor to categories
is faithful but not full onto all endpoint-preserving interval
functors. The vertex sequence $(0,1,0,1,2,3)$ defines a functor
$I_{+-+-+}\to I_{+-+}$, since its five arrow images are respectively
$0\to1,0\to1,0\to1,2\to1,2\to3$. It preserves endpoints and is not
nondecreasing. This checks the precise restriction in the written proof.

\textbf{Failure of filteredness.} By the allocation description,
$J=\mathrm{Ty}^{\op}$ is the category of finite linearly ordered
sets colored by $+,-$ and monotone color-preserving maps. The two maps
$+\rightrightarrows +-+$ selecting its first and last positions have
no common equalizing postcomposition. Such a postcomposition would
identify the two positive positions and hence also the intervening
negative position by monotonicity, contradicting color preservation.

\textbf{An explicit witness against total asphericity.} Put $p=++$
and $S(t)=J(p,t)$. The category $E=\operatorname{El}_J(S)=p\downarrow J$
has initial object $1_p$. An object is a pair of positive positions
$i\leq j$ in a word $t$. Let $U$ be the cosieve generated by the
collapse $p\to+$, and $V$ the cosieve generated by the map
$p\to+-+$ selecting the outer positions. They have the following
exact descriptions: $U$ consists of equal pairs $i=j$; $V$ consists
of pairs with a negative position strictly between them. They are
nonempty, connected, disjoint cosieves. For example an image of a
negative position strictly between two positive ones must remain
strictly between them, because colors prevent equality.

On $E$ use the two-element action outside $U\cup V$ and the singleton
action on $U\cup V$, with identity or collapse structure maps. Inducing
this action to $J$ gives the explicit covariant set functor $L$: a pair
$i\leq j$ has one copy if $i=j$ or a negative position lies between
them, and two copies otherwise. A monotone colored map sends the pair
to its image and retains its copy label unless the target has one
copy. The cosieve property proves functoriality. Its element category
is two copies of $E$ glued on $U\sqcup V$. A chain cannot leave a
cosieve after entering it, so this is also a pushout of nerves.
Both copies of $BE$ are contractible; $BU$ and $BV$ are connected.
Mayer--Vietoris consequently gives
\[
 \widetilde H_0(B\operatorname{El}L;\Z)=0,
 \qquad H_1(B\operatorname{El}L;\Z)=\Z.
\]
For a direct test of the definition, adjoin to $J$ a new source $*$
with $\Hom(*,t)=L(t)$, no arrows back, and its identity endomorphism.
Call the resulting category $J'$ and the inclusion $j:J\to J'$.
The comma $*\downarrow j$ is $\operatorname{El}L$ and is connected;
the comma at an old object has an initial object. Therefore the
ordinary presheaf left Kan extension satisfies $j_!1_J=1_{J'}$.
The nerve of this collage gives the homotopy cofiber of
$B\operatorname{El}L\to BJ$. Since $J$ has the initial empty word,
$BJ$ is contractible and $H_2(BJ';\Z)=\Z$. Thus the element category
of this Kan extension is not weakly contractible. This refutes total
$W_\infty$-asphericity under exactly the definition in
\textsc{PRETOT}, rather than inferring it from failure of filteredness.

\textbf{The actual sequential colimit.} Pullback reverses type maps,
so a projective sequence acts covariantly on paths with transitions
$K_{n-1}\to K_n$. Nerves commute with filtered categorical colimits:
an $r$-simplex and each equality between two such simplices are finite
lists of arrows and composition equations. The simplicial replacement
map from the filtered homotopy colimit to this colimit is a weak
equivalence. One can check this directly on homotopy groups: a
triangulated sphere, disk, and any homotopy have finitely many simplices;
their data and equations have a common stage, and a finite set of
stages and arrows has a cone at a later stage. These cones give the
required lifts and equality witnesses. Realization yields the stated
homotopy equivalence, independently of any cofinality of the sequence
inside all types.

The independent program \texttt{f\_type\_categories.py} checks all
$961$ pairs of words of length at most four, including $1471$ type
maps and $10771$ transformation tests, against the interval definition.
It also checks the complete colored-order category on words of length
at most three ($15$ objects, $209$ arrows) and the above action; its
element category has $30$ objects, $499$ arrows, $94404$ composable
triples and first mod-two Betti number one. These are checks on fixed
finite instances; the universal assertions use the arguments above.
\end{proof}

\begin{proposition}[The first proof of C.20, with its missing prerequisite supplied]
\label{F:C20-paths}
Let $(M,\Sigma)$ be a small relative category with all identities and
composites in its marking. Let $t:\N^{\op}\to\mathrm{Ty}$ be an actual
sequence such that every type receives a map from some $t_n$.
Write $K_n=F_{t_n}(a,b)$ for paths with marked negative edges, arbitrary
interior vertical arrows, and identity endpoint components. For every
small category $X$, including infinite and empty ones, effect of words
gives a natural bijection
\[
 \colim_n\pi_0\operatorname{Fun}(X,K_n)
 \cong
 \Hom_{M^X[\Sigma_X^{-1}]}(a_X,b_X).
\]
Here $\Sigma_X$ consists of the pointwise marked transformations.
The comparison respects constant paths, concatenation, and restriction
in $X$.
\end{proposition}
\begin{proof}
\textbf{The exponential-law step.} A path of type $\tau$ in $M^X$ is
exactly a functor $X\times I_\tau\to M$ with the stated endpoint and
marking conditions. Its ladders have exactly the endpoint-fixed
components in the definition. Hence
$F_\tau^{M^X}(a_X,b_X)=\operatorname{Fun}(X,F_\tau^M(a,b))$
canonically, on objects and arrows, and compatibly with type pullbacks.

\textbf{The component universal property.} Compose forward arrows and
inverses of backward marked arrows in the ordinary localization.
Type pullback does not change this word's effect. A positive ladder
square preserves the composite equation; a negative square does so
after cancelling its two invertible horizontal edges. No interior
vertical edge is inverted. Conversely components of all-type paths
form a category under concatenation. Selecting the outer vertices
compares two successive forward edges with their composite; collapse
compares an inserted identity. For $w:x\to y$ in the marking, compare
the backtrack $x\xrightarrow{w}y\xleftarrow{w}x$ to the constant
path at $x$ by the ladder with middle component $w:x\to y$ and
identity endpoint components. For the other backtrack use the ladder
to the constant path at $y$, again with middle component $w$.
Thus marked arrows become invertible. Evaluation in any category
inverting $\Sigma_X$ respects all these equations and gives its unique
factorization. This proves the localization universal property and
identifies the specified effect map.

\textbf{Surjectivity and injectivity at actual stages.} Each word has
one finite type, so pulling it back to a $t_n$ proves surjectivity.
An equality of components is a finite zigzag of total-path arrows.
Concatenate all its types, and choose a stage mapping to that word,
increasing the stage to dominate both original endpoints. Pull every
ladder back there. A type-change arrow may compare two pullbacks of
one path along different maps from this stage to the same type.
Lemma~\ref{F:type-contraction} connects these maps by transformations
uniformly in the path; composing and applying $\operatorname{Fun}(X,-)$
gives the required ladder zigzag even for infinite $X$. The same
argument at both endpoints compares the choices with the actual
transition maps of the original sequence. The equality therefore holds
at an actual common stage. Every step uses restriction, word effect,
or concatenation, proving the claimed naturalities and units.

\textbf{The exact marking hypothesis.} Pullback along
$-\to\varnothing$ requires marked identities; pullback along
$-\to--$ requires closure under composition. These conditions suffice
because a negative edge maps to a negative string or an identity.
For an arbitrary class of arrows, replacing it by its generated wide
subcategory preserves the ordinary localization by its universal
property and makes the type diagram defined. This confirms all steps
of the supplied first proof of C.20; its alternative proof via C.3
is not needed for this assertion.
\end{proof}
